\documentclass[../main.tex]{subfiles}
\usepackage[utf8]{inputenc}
\usepackage[T1]{fontenc}
\usepackage[indonesian]{babel}

\begin{document}

\chapter{Kriptografi Modern Simetris: Block Cipher}

\begin{learningoutcome}
Setelah mempelajari bab ini, mahasiswa diharapkan mampu:
\begin{itemize}
    \item Menjelaskan konsep dasar \textit{Block Cipher} dan perbedaannya dengan \textit{Stream Cipher} (Sub-CPMK 1.2).
    \item Menganalisis berbagai mode operasi \textit{Block Cipher} (ECB, CBC, CFB, OFB, CTR) serta kelebihan dan kelemahannya (Sub-CPMK 1.2).
    \item Memahami prinsip \textit{Confusion} dan \textit{Diffusion} dalam desain algoritma kriptografi (Sub-CPMK 1.2).
    \item Menjelaskan struktur \textit{Iterated Cipher} dan mekanisme Jaringan Feistel (Sub-CPMK 1.2).
    \item Memahami proses pembangkitan kunci putaran (\textit{key scheduling}) dan pengaruh ukuran blok serta kunci terhadap keamanan (Sub-CPMK 1.2).
\end{itemize}
\end{learningoutcome}

\section{Konsep Dasar Block Cipher}
Berbeda dengan \textit{stream cipher} yang memproses data bit-per-bit, \textit{block cipher} membagi plainteks menjadi blok-blok bit dengan panjang yang sama.

\subsection{Karakteristik Utama}
\begin{itemize}
    \item \textbf{Ukuran Blok}: Plainteks dibagi menjadi blok berukuran $n$ bit. Ukuran blok yang umum digunakan adalah 64 bit, 128 bit, atau 256 bit.
    \item \textbf{Panjang Output}: Panjang blok cipherteks selalu sama dengan panjang blok plainteks.
    \item \textbf{Operasi}: Enkripsi dilakukan pada setiap blok plainteks menggunakan kunci $K$.
\end{itemize}

Secara matematis, jika $P$ adalah blok plainteks dan $C$ adalah blok cipherteks, maka:
\[ C = E_K(P) \quad \text{dan} \quad P = D_K(C) \]
di mana $E$ adalah fungsi enkripsi dan $D$ adalah fungsi dekripsi.

\subsection{Contoh Sederhana Block Cipher}
Sebagai ilustrasi, misalkan sebuah fungsi enkripsi sederhana yang melakukan XOR antara plainteks $P$ dengan kunci $K$, kemudian menggeser bit-bitnya satu posisi ke kiri secara \textit{wrapping}:
\[ E_K(P) = (P \oplus K) \ll 1 \]
Maka fungsi dekripsinya adalah kebalikan dari proses tersebut:
\[ D_K(C) = (C \gg 1) \oplus K \]

\section{Mode Operasi Block Cipher}
Karena \textit{block cipher} hanya dapat mengenkripsi satu blok data, diperlukan \textbf{mode operasi} untuk mengenkripsi pesan yang panjangnya lebih dari satu blok.

\subsection{Electronic Code Book (ECB)}
Mode paling sederhana di mana setiap blok plainteks $P_i$ dienkripsi secara independen:
\[ C_i = E_K(P_i) \]
\textbf{Kelemahan}: Blok plainteks yang sama akan selalu menghasilkan blok cipherteks yang sama. Hal ini menyebabkan pola data asli tetap terlihat pada cipherteks, sehingga rentan terhadap analisis statistik.

\subsection{Cipher Block Chaining (CBC)}
Untuk mengatasi kelemahan ECB, CBC memperkenalkan ketergantungan antar blok. Setiap blok plainteks di-XOR dengan blok cipherteks sebelumnya sebelum dienkripsi:
\[ C_0 = IV \quad (\text{Initialization Vector}) \]
\[ C_i = E_K(P_i \oplus C_{i-1}) \]
Dekripsinya adalah:
\[ P_i = D_K(C_i) \oplus C_{i-1} \]
CBC memastikan bahwa blok plainteks yang sama menghasilkan cipherteks yang berbeda tergantung posisinya.

\subsection{Cipher Feedback (CFB) dan Output Feedback (OFB)}
Kedua mode ini mengubah \textit{block cipher} agar bekerja menyerupai \textit{stream cipher}:
\begin{itemize}
    \item \textbf{CFB}: Hasil enkripsi dari blok sebelumnya di-XOR dengan plainteks saat ini. Cocok untuk data yang belum mencapai ukuran satu blok penuh.
    \item \textbf{OFB}: Mirip dengan CFB, tetapi yang di-XOR dengan plainteks adalah keluaran dari generator kunci internal (tidak bergantung pada cipherteks). Ini mencegah penjalaran kesalahan (\textit{error propagation}).
\end{itemize}

\subsection{Counter Mode (CTR)}
Mode CTR mengenkripsi sebuah nilai \textit{counter} yang terus bertambah, kemudian hasil enkripsinya di-XOR dengan plainteks:
\[ C_i = P_i \oplus E_K(\text{counter}_i) \]
\textbf{Kelebihan}: Mendukung pemrosesan paralel karena setiap blok tidak bergantung pada blok sebelumnya.

\section{Prinsip Confusion dan Diffusion}
Claude Shannon memperkenalkan dua prinsip utama dalam desain cipher untuk menangkal kriptanalisis statistik:

\subsection{Confusion (Konfusi)}
\textit{Confusion} bertujuan untuk membuat hubungan antara kunci dan cipherteks menjadi serumit mungkin. Hal ini biasanya direalisasikan melalui teknik \textbf{substitusi} non-linear, seperti penggunaan \textbf{S-Box} (Substitution Box) yang memetakan input ke output berdasarkan tabel pencarian.

\subsection{Diffusion (Difusi)}
\textit{Diffusion} bertujuan untuk menyebarkan pengaruh satu bit plainteks ke banyak bit cipherteks. Artinya, perubahan satu bit pada plainteks harus menyebabkan perubahan signifikan (sekitar setengah) pada bit-bit cipherteks. Hal ini direalisasikan melalui teknik \textbf{permutasi} atau transposisi (misal: $P$-box atau \textit{ShiftRows} pada AES).

\section{Iterated Cipher dan Jaringan Feistel}
Untuk mencapai keamanan yang tinggi, \textit{block cipher} modern menggunakan struktur \textbf{Iterated Cipher}, di mana fungsi transformasi (fungsi putaran) diterapkan berulang kali selama $r$ putaran.

\subsection{Jaringan Feistel}
Struktur yang sangat populer (digunakan pada DES) adalah \textbf{Jaringan Feistel}. Dalam struktur ini, blok data dibagi menjadi dua bagian, Kiri ($L$) dan Kanan ($R$):
\begin{align*}
L_i &= R_{i-1} \\
R_i &= L_{i-1} \oplus f(R_{i-1}, K_i)
\end{align*}
Sifat utama Jaringan Feistel adalah \textbf{reversible}. Proses dekripsi dilakukan dengan urutan putaran yang terbalik menggunakan kunci yang sama, tanpa perlu membuat fungsi dekripsi $f^{-1}$ yang terpisah.

\subsection{Pembangkitan Kunci Putaran (Key Scheduling)}
Setiap putaran dalam \textit{iterated cipher} menggunakan kunci internal yang berbeda, disebut \textbf{subkey} atau \textit{round key}. Kunci-kunci ini dibangkitkan dari kunci eksternal yang diberikan pengguna melalui proses \textit{key expansion}.

\subsection{Ukuran Blok dan Kunci}
\begin{itemize}
    \item \textbf{Ukuran Blok}: Semakin besar blok, semakin kuat efek difusinya (misal: AES menggunakan 128 bit).
    \item \textbf{Ukuran Kunci}: Semakin panjang kunci, semakin tahan cipher terhadap serangan \textit{brute force} (misal: AES-256 lebih kuat dari AES-128).
\end{itemize}

\begin{obeactivity}
\textbf{Aktivitas 6.1: Analisis Mode ECB vs CBC}
1. Gunakan alat bantu enkripsi (seperti CyberChef) untuk mengenkripsi sebuah gambar sederhana (misal: gambar logo kecil) menggunakan AES mode ECB dan AES mode CBC.
2. Amati hasil cipherteks yang dihasilkan. Apakah gambar asli masih terlihat pada mode ECB?
3. Diskusikan mengapa mode CBC lebih aman untuk data yang memiliki pola berulang.
\end{obeactivity}

\begin{obereflection}
\textbf{Latihan 6.1}
1. Jika sebuah blok cipher memiliki ukuran 64 bit dan kita menggunakan mode ECB, apa yang terjadi jika dua blok plainteks identik dikirimkan?
2. Jelaskan perbedaan mendasar antara peran S-Box dan P-Box dalam mencapai \textit{confusion} dan \textit{diffusion}.
3. Mengapa Jaringan Feistel memudahkan implementasi dekripsi?
\end{obereflection}

\begin{obeassessment}
\textbf{Kuis Bab 6}
1. Berikan contoh skenario di mana mode CTR lebih menguntungkan dibandingkan mode CBC.
2. Apa yang dimaksud dengan \textit{error propagation} pada mode CBC dan bagaimana mode OFB mengatasinya?
3. Mengapa penggunaan IV (\textit{Initialization Vector}) sangat penting pada mode CBC?
\end{obeassessment}

\begin{competencychecklist}
\begin{tabular}{|l|c|}
\hline
Indikator Kompetensi & Tercapai \\
\hline
Saya dapat menjelaskan perbedaan \textit{block cipher} dan \textit{stream cipher} & [ ] \\
\hline
Saya dapat membedakan mode operasi ECB, CBC, CFB, OFB, dan CTR & [ ] \\
\hline
Saya dapat menjelaskan prinsip \textit{confusion} dan \textit{diffusion} & [ ] \\
\hline
Saya dapat menguraikan mekanisme kerja Jaringan Feistel & [ ] \\
\hline
Saya dapat menganalisis pengaruh panjang kunci terhadap keamanan sistem & [ ] \\
\hline
\end{tabular}
\end{center}
\end{competencychecklist}

\end{document}
